% Bagian: first-order-logic
% Bab: completeness
% Subbagian: maximally-consistent-sets

% Definisi himpunan konsisten maksimal. Sifat-sifat himpunan konsisten
% maksimal yang diperlukan untuk kelengkapan dibuktikan dalam
% maximally-consistent-sets.tex pada bab tentang sistem bukti yang digunakan.

\documentclass[../../../include/open-logic-section]{subfiles}

\begin{document}

\olfileid[id]{fol}{com}{mcs}
\olsection{Himpunan \usetoken{P}{sentence} yang Konsisten Maksimal}

\begin{defn}[Himpunan konsisten maksimal]
\ollabel{mcs}
Suatu himpunan~$\Gamma$ dari !!{sentence} bersifat \emph{konsisten maksimal}
jika dan hanya jika
\begin{enumerate}
\item $\Gamma$ konsisten, dan
\item jika $\Gamma \subsetneq \Gamma'$, maka $\Gamma'$ tak konsisten.
\end{enumerate}
\end{defn}

\begin{explain}
Definisi lain yang ekuivalen dengan definisi di atas ialah: suatu himpunan
kalimat~$\Gamma$ bersifat \emph{konsisten maksimal} jika dan hanya jika
\begin{enumerate}
\item $\Gamma$ konsisten, dan
\item jika $\Gamma \cup \{ !A \}$ konsisten, maka $!A \in \Gamma$.
\end{enumerate}
Dengan kata lain, !!{sentence} yang belum berada dalam~$\Gamma$ tidak dapat
ditambahkan kepada himpunan konsisten maksimal~$\Gamma$ tanpa membuat
himpunan lebih besar yang dihasilkan menjadi tak konsisten.
\end{explain}

\begin{explain}
Himpunan konsisten maksimal penting dalam pembuktian kelengkapan karena kita
dapat menjamin bahwa setiap himpunan konsisten dari !!{sentence}~$\Gamma$
termuat dalam suatu himpunan konsisten maksimal~$\Gamma^*$, dan suatu himpunan
konsisten maksimal memuat, untuk setiap !!{sentence}~$!A$, $!A$ atau
negasinya~$\lnot !A$. Hal ini berlaku khususnya bagi !!{sentence} atomik;
karena itu, dari suatu himpunan konsisten maksimal dalam bahasa yang diperluas
secara tepat dengan !!{constant}, kita dapat membangun !!a{structure} dengan
interpretasi !!{predicate} didefinisikan menurut !!{sentence} atomik mana
yang berada dalam~$\Gamma^*$. Selanjutnya dapat ditunjukkan bahwa
!!{structure} ini membuat semua !!{sentence} dalam~$\Gamma^*$ (dan karena itu
juga semua kalimat dalam~$\Gamma$) benar. Pembuktian fakta terakhir ini
mensyaratkan bahwa $\lnot !A \in \Gamma^*$ jika dan hanya jika
$!A \notin \Gamma^*$, bahwa $(!A \lor !B) \in \Gamma^*$ jika dan hanya
jika $!A \in \Gamma^*$ atau $!B \in \Gamma^*$, dan seterusnya.
\end{explain}

\begin{prop}
\ollabel{prop:mcs}
Andaikan $\Gamma$ konsisten maksimal. Maka:
\begin{enumerate}
\item \ollabel{prop:mcs-prov-in} Jika $\Gamma \Proves !A$, maka $!A \in
  \Gamma$.

\item \ollabel{prop:mcs-either-or} Untuk setiap $!A$, berlaku $!A \in
  \Gamma$ atau $\lnot !A \in \Gamma$.

\tagitem{prvAnd}{\ollabel{prop:mcs-and} $(!A \land !B) \in \Gamma$
  jika dan hanya jika $!A \in \Gamma$ dan $!B \in \Gamma$.}{}

\tagitem{prvOr}{\ollabel{prop:mcs-or} $(!A \lor !B) \in \Gamma$ jika dan
  hanya jika $!A \in \Gamma$ atau $!B \in \Gamma$.}{}

\tagitem{prvIf}{\ollabel{prop:mcs-if} $(!A \lif !B) \in \Gamma$ jika dan
  hanya jika $!A \notin \Gamma$ atau $!B \in \Gamma$.}{}
\end{enumerate}
\end{prop}

\begin{proof}
Untuk seluruh pembahasan berikut, andaikan bahwa $\Gamma$ konsisten maksimal.
\begin{enumerate}
\item Jika $\Gamma \Proves !A$, maka $!A \in \Gamma$.

Andaikan $\Gamma \Proves !A$. Demi memperoleh kontradiksi, andaikan bahwa
$!A \notin \Gamma$. Karena $\Gamma$ konsisten maksimal, $\Gamma \cup
\{!A\}$ tak konsisten, sehingga $\Gamma \cup \{!A\} \Proves \lfalse$.
Menurut
\tagrefs{prfSC/{fol:seq:prv:prop:provability-contr},prfND/{fol:ntd:prv:prop:provability-contr}},
$\Gamma$ tak konsisten. Hal ini bertentangan dengan asumsi bahwa $\Gamma$
konsisten. Jadi, tidak mungkin $!A \notin \Gamma$, sehingga
$!A \in \Gamma$.

\item Untuk setiap $!A$, berlaku $!A \in \Gamma$ atau
$\lnot !A \in \Gamma$.

Demi memperoleh kontradiksi, andaikan bahwa untuk suatu~$!A$ berlaku
$!A \notin \Gamma$ dan $\lnot !A \notin \Gamma$. Karena $\Gamma$
konsisten maksimal, $\Gamma \cup \{!A\}$ dan
$\Gamma \cup \{\lnot !A\}$ keduanya tak konsisten, sehingga
$\Gamma \cup \{!A\} \Proves \lfalse$ dan
$\Gamma \cup \{\lnot !A\} \Proves \lfalse$. Menurut
\tagrefs{prfSC/{fol:seq:prv:prop:provability-exhaustive},prfND/{fol:ntd:prv:prop:provability-exhaustive}},
$\Gamma$ tak konsisten, suatu kontradiksi. Jadi, !!a{sentence}~$!A$ semacam
itu tidak mungkin ada dan, untuk setiap $!A$, berlaku $!A \in \Gamma$ atau
$\lnot !A \in \Gamma$.

\tagitem{defAnd}{}{%
\iftag{probAnd}{Latihan.}{%
$(!A \land !B) \in \Gamma$ jika dan hanya jika $!A \in \Gamma$ dan
$!B \in \Gamma$:

Untuk arah maju, andaikan $(!A \land !B) \in \Gamma$. Maka
$\Gamma \Proves !A \land !B$. Menurut
\tagrefs{prfSC/{fol:seq:prv:prop:provability-land-left},prfND/{fol:ntd:prv:prop:provability-land-left}},
$\Gamma \Proves !A$ dan $\Gamma \Proves !B$. Menurut
\olref{prop:mcs-prov-in}, $!A \in \Gamma$ dan $!B \in \Gamma$, seperti
yang diperlukan.

Untuk arah sebaliknya, misalkan $!A \in \Gamma$ dan $!B \in \Gamma$.
Maka $\Gamma \Proves !A$ dan $\Gamma \Proves !B$. Menurut
\tagrefs{prfSC/{fol:seq:prv:prop:provability-land-right},prfND/{fol:ntd:prv:prop:provability-land-right}},
$\Gamma \Proves !A \land !B$. Menurut \olref{prop:mcs-prov-in},
$(!A \land !B) \in \Gamma$.}}

\tagitem{defOr}{}{%
\iftag{probOr}{Latihan.}{%
$(!A \lor !B) \in \Gamma$ jika dan hanya jika $!A \in \Gamma$ atau
$!B \in \Gamma$.

Untuk kontrapositif arah maju, andaikan bahwa $!A \notin \Gamma$ dan
$!B \notin \Gamma$. Kita hendak menunjukkan bahwa
$(!A \lor !B) \notin \Gamma$. Karena $\Gamma$ konsisten maksimal,
$\Gamma \cup \{!A\} \Proves \lfalse$ dan
$\Gamma \cup \{!B\} \Proves \lfalse$. Menurut
\tagrefs{prfSC/{fol:seq:prv:prop:provability-lor-left},prfND/{fol:ntd:prv:prop:provability-lor-left}},
$\Gamma \cup \{(!A \lor !B)\}$ tak konsisten. Jadi,
$(!A \lor !B) \notin \Gamma$, seperti yang diperlukan.

Untuk arah sebaliknya, andaikan bahwa $!A \in \Gamma$ atau
$!B \in \Gamma$. Maka $\Gamma \Proves !A$ atau $\Gamma \Proves !B$.
Menurut
\tagrefs{prfSC/{fol:seq:prv:prop:provability-lor-right},prfND/{fol:ntd:prv:prop:provability-lor-right}},
$\Gamma \Proves !A \lor !B$. Menurut \olref{prop:mcs-prov-in},
$(!A \lor !B) \in \Gamma$, seperti yang diperlukan.}}

\tagitem{defIf}{}{%
\iftag{probIf}{Latihan.}{%
$(!A \lif !B) \in \Gamma$ jika dan hanya jika $!A \notin \Gamma$ atau
$!B \in \Gamma$:

Untuk arah maju, misalkan $(!A \lif !B) \in \Gamma$, lalu, demi memperoleh
kontradiksi, andaikan bahwa $!A \in \Gamma$ dan $!B \notin \Gamma$.
Berdasarkan asumsi-asumsi ini, $\Gamma \Proves !A \lif !B$ dan
$\Gamma \Proves !A$. Menurut
\tagrefs{prfSC/{fol:seq:prv:prop:provability-mp},prfND/{fol:ntd:prv:prop:provability-mp}},
$\Gamma \Proves !B$. Namun, menurut \olref{prop:mcs-prov-in}, hal ini
berarti $!B \in \Gamma$, bertentangan dengan asumsi bahwa
$!B \notin \Gamma$.

Untuk arah sebaliknya, pertama-tama tinjau kasus ketika
$!A \notin \Gamma$. Menurut \olref{prop:mcs-either-or},
$\lnot !A \in \Gamma$, sehingga $\Gamma \Proves \lnot !A$. Menurut
\tagrefs{prfSC/{fol:seq:prv:prop:provability-lif},prfND/{fol:ntd:prv:prop:provability-lif}},
$\Gamma \Proves !A \lif !B$. Sekali lagi, menurut
\olref{prop:mcs-prov-in}, kita memperoleh $(!A \lif !B) \in \Gamma$,
seperti yang diperlukan.

Sekarang tinjau kasus ketika $!B \in \Gamma$. Maka $\Gamma \Proves !B$,
dan menurut
\tagrefs{prfSC/{fol:seq:prv:prop:provability-lif},prfND/{fol:ntd:prv:prop:provability-lif}},
$\Gamma \Proves !A \lif !B$. Menurut \olref{prop:mcs-prov-in},
$(!A \lif !B) \in \Gamma$.}}
\end{enumerate}
\end{proof}

\begin{prob}
Lengkapi pembuktian \olref[fol][com][mcs]{prop:mcs}.
\end{prob}

\end{document}
